\part{Institutions, Power, and Distribution}
\label{part:VI}

Part VI analyzes classifications, property, contracts, exit, and public authority as distinctions that alter which continuations are accessible and to whom.

\begin{housing}Tenure, deposits, notice periods, and a tenant's executable capacity to leave are placed in one model.\end{housing}

\chapter{Distinctions That Allocate Futures}\label{ch:26}
Examines property categories, employment classifications, citizenship, credit status, and professional credentials as distinctions that alter accessible continuations. Connects the theory of distinction and continuation to institutional economics.
\scaffold

\chapter{Property as Structured Authority}\label{ch:27}
Analyzes property through permissions to use, exclude, modify, transfer, encumber, and abandon. Distinguishes ownership from the practical capacity to exercise those permissions and from the burdens imposed on others.
\scaffold

\chapter{Contracts and the Evolution of Dependence}\label{ch:28}
Studies contracts as trajectory-shaping arrangements. Examines how repeated agreements alter bargaining positions, switching costs, and the conditions of later consent, and analyzes cumulative dependence under continuing intervention.
\scaffold

\chapter{Distribution of Exit and Refusal}\label{ch:29}
Treats inequality as including unequal capacities to leave employment, housing, platforms, financial relationships, and harmful environments. Distinguishes nominal exit rights from executable exit paths, and states when an intervention can raise surplus while shrinking a participant's safe exit set.
\chspec{\Pref{P08}, \Pref{P28}.}
\scaffold

\chapter{Public Authority and the Admissibility of Constraints}\label{ch:30}
Addresses democratic authorization, contestability, rights, and institutional revision. Examines who establishes continuation constraints, how conflicts among them are adjudicated, and what prevents preservation language from legitimating domination.
\draftnote{Candidate question: what must happen before a repair obligation may legitimately be discharged? A waiting period is an enforceable procedural condition but elapsed time cannot establish successful repair. Closure should require procedural authorization and specified evidence, with a route for reopening a failed repair. Grouping reports by location suppresses duplicate influence but does not establish independent witnesses, and a minimum cluster size does not by itself protect against all differencing attacks. Present as a proposed institutional design, not evidence of deployment.}
\scaffold
